\section{Introduction}
\label{sec:intro}

\textit{Arabidopsis thaliana} grows in a light environment that switches off
every day. During the light phase the plant assimilates carbon and stores part
of it as transient leaf starch. During the night this reserve is the only carbon
source for respiration, maintenance and growth. Measurements show that the
reserve is degraded at an almost constant rate and is nearly exhausted at dawn
\cite{Graf2010,Smith2007}. The plant therefore behaves as if it rations a known
amount of carbon over a night of known length \cite{Scialdone2013}. When the
photoperiod changes, both the amount and the length change together, and the
rationing problem changes with them \cite{Sulpice2014,Stitt2012}. This makes
photoperiod a useful test case for a simulation model, because a single
environmental variable moves the light-phase carbon gain and the dark-phase
reserve use at the same time.

The problem is difficult to simulate because it couples processes on very
different time scales. Stomatal opening and leaf assimilation respond within
minutes to light, temperature and humidity. Starch and sucrose pools turn over
within hours. Nutrient status changes over days and biomass accumulates over
weeks. The carbon available for growth at any instant is not a property of any
single equation. It is the result of the environment, the gas-exchange
machinery and the internal carbohydrate and nutrient state acting together
\cite{Farquhar1980,ThornleyJohnson1990}. A model that describes only
photosynthesis, or only biomass accumulation, cannot represent this. For a
simulation study the main task is therefore to build a stable numerical
workflow in which these subsystems exchange fluxes consistently over the
day--night cycle.

Existing models resolve parts of this chain, but rarely the whole chain at the
resolution the problem needs. Crop and field-scale simulators integrate growth,
water and nitrogen over seasons, but their standard configurations prescribe the
carbon input and advance the state on a daily step
\cite{deWit2019,Holzworth2014,Jones2003}. Leaf gas-exchange models resolve the
process that a prescribed carbon input replaces, but are formulated and tested as
steady-state systems at a fixed carbohydrate status, so they stop at the leaf
\cite{Farquhar1980,Jarvis1976,Damour2010}. Few published models compute the
supply--demand operating point at every step and then carry that carbon input
through a diel starch and sucrose economy into biomass. Section~\ref{sec:related}
develops this gap in detail.

This paper presents such a model and reports its evaluation. The model is a
whole-plant dynamic simulator with eleven state variables, organised as six
submodels with declared inputs, outputs and parameters
(Figure~\ref{fig:overview}). S1 maps clock time to
the environmental drivers. S2 computes stomatal conductance from a
multiplicative Jarvis formulation, which acts as the supply side of gas
exchange. S3 computes net assimilation as a function of intercellular CO$_2$
from the Farquhar C3 biochemistry, which acts as the demand side. S4 closes
supply against demand. S5 holds the diel carbohydrate economy, the starch and
sucrose pools and the nitrogen and phosphorus uptake and costs. S6 converts
sucrose into leaf and root biomass. Submodels S2--S4 form a fast algebraic
subsystem, and S5--S6 form the slow differential subsystem that is integrated
in time.

The methodological content sits in the runtime solution and the evaluation
design, and this is what we expect to transfer to other coupled biophysical
simulators. The gas-exchange operating point is solved algebraically at runtime:
substituting Fick's law into the branch-wise Farquhar expression gives a
quadratic in the intercellular CO$_2$ concentration, solved branch by branch
with an explicit admissibility rule for the roots, so no inner iteration is
required inside the derivative evaluation. The algebraic subsystem is then
eliminated exactly inside each derivative call, and the remaining differential
system is integrated with an explicit adaptive method. The implementation is
verified before it is compared with data, and model behaviour is analysed with a
variance-based global sensitivity analysis over declared parameter ranges. The
contributions are listed at the end of this section.

The evaluation design keeps calibration and prediction separate. Four
coefficients are fitted and shared across all photoperiods, in place of a legacy
per-photoperiod calibration that left no degrees of freedom
(Section~\ref{sec:protocol}). Predictive accuracy is then measured by holding out
whole photoperiods: fitting on three light regimes and predicting the fourth.
Under this design the model reaches a held-out mean absolute percentage error of
4.36\,\% in relative growth rate, when one growth coefficient is allowed to
depend linearly on photoperiod (Figure~\ref{fig:crossval}). We also report the
photoperiod range over which the model remains usable, together with the
mechanism that breaks first.

The contributions of the paper are the following.
\begin{enumerate}
\item A coupled whole-plant simulation model in which Jarvis stomatal
  conductance and Farquhar C3 photosynthesis produce the carbon input at every
  step, and that input is propagated through diel starch and sucrose dynamics,
  nutrient uptake and costs, and biomass growth (Section~\ref{sec:model}).
\item An analytic runtime solution of the supply--demand closure as a quadratic
  in the intercellular CO$_2$ concentration, embedded in a fast/slow numerical
  workflow (Sections~\ref{sec:model} and~\ref{sec:implementation}).
\item A verification protocol for the implementation that is independent of any
  comparison with data, covering the algebraic closure, the time integration
  and the carbohydrate bookkeeping (Section~\ref{sec:implementation}).
\item An evaluation with four shared fitted parameters instead of a
  per-photoperiod calibration. The design uses leave-one-photoperiod-out
  cross-validation, an ablation of the carbon-input pathway, and a global
  sensitivity analysis over declared ranges
  (Sections~\ref{sec:protocol} and~\ref{sec:results}).
\item A measured validity envelope. We report the photoperiod range in which
  the diel starch balance closes, the photoperiod at which it stops closing,
  and the bound on the simulated growth response to day length
  (Sections~\ref{sec:results} and~\ref{sec:limitations}).
\end{enumerate}

The rest of the paper is organised as follows.
Section~\ref{sec:related} reviews mechanistic growth models, gas-exchange
formulations, diel carbohydrate modelling and simulation evaluation practice,
and states the gap that this work addresses. Section~\ref{sec:model} gives the
model formulation as six submodels with their states, inputs and parameters.
Section~\ref{sec:implementation} describes the numerical implementation and the
verification results. Section~\ref{sec:protocol} defines the photoperiod
scenarios, the calibration and cross-validation design, the ablation variants
and the sensitivity protocol. Section~\ref{sec:results} reports the results,
Section~\ref{sec:discussion} interprets them, Section~\ref{sec:limitations}
states what cannot be concluded, and Section~\ref{sec:conclusion} concludes.

% ============================================================
\section{Related Work}
\label{sec:related}

This section positions the model against four bodies of work: mechanistic
whole-plant and crop simulation, leaf gas-exchange modelling, diel carbohydrate
turnover in \textit{Arabidopsis thaliana}, and evaluation practice in
simulation studies. The aim is to locate the contribution, not to survey each
field.

\subsection{Mechanistic plant growth simulation models}

Mechanistic growth models link physiological processes, microclimate and
biomass change over time, at scales from the leaf to the field and from minutes
to whole seasons. Crop simulation platforms such as APSIM, DSSAT and WOFOST are
mature tools for growth, water and nitrogen dynamics at field scale
\cite{Holzworth2014,Jones2003,deWit2019}, and process-based formulations follow
seasonal climate variation better than empirical ones
\cite{BuckleyPaules2025,Mamassi2024}. Other families address the same problem
differently: functional--structural models combine architecture with
physiology \cite{Louarn2020,Sievnen2014,DeReffye2020}, soil--plant water
transport models describe flow from soil through xylem \cite{Janott2011},
optimality models assume allocation maximises a criterion \cite{Russo2022}, and
continuous formulations describe biomass or vegetation dynamics as differential
equations \cite{Meron2019,ThornleyJohnson1990,Marcelis1996,Feller2015,Thornley1972}.

Two properties of this literature matter here. Modular organisation is common
and useful, because it lets each component be tested on its own
\cite{Sievnen2014,Zeigler2000}. In the configurations most often used at field
scale, the carbon input is prescribed rather than computed, through
light-use-efficiency or radiation-use coefficients with a daily step
\cite{Jones2003,deWit2019,Holzworth2014}. That default separates the carbon
source from the environment and from the internal state of the plant, so the
model cannot represent the feedback between stomatal behaviour, intercellular
CO$_2$, carbohydrate status and growth. This describes the usual configuration,
not the capability of these platforms.

\subsection{Photosynthesis and stomatal conductance modelling}

For C3 plants the standard biochemical model is the
Farquhar--von Caemmerer--Berry formulation, which relates net assimilation to
intercellular CO$_2$ through Rubisco-limited carboxylation and RuBP regeneration
\cite{Farquhar1980,vonCaemmerer2000}; its parameters and their temperature
dependence are well documented \cite{Sharkey2007,Bernacchi2001,Dubois2007}.
Assimilation also depends on the diffusive supply of CO$_2$ set by stomatal
conductance \cite{Lawson2014}. The main conductance formulations are the
empirical Jarvis model, the Ball--Berry family and the Medlyn optimality model
\cite{Jarvis1976,Ball1987,Medlyn2011}. The Jarvis form expresses conductance as
a product of response functions to irradiance, vapour-pressure deficit,
temperature, CO$_2$ and plant water status \cite{Jarvis1976,YU2017,Yu2021}, and
is convenient in simulation because time-varying drivers enter it directly
\cite{Damour2010,Buckley2017}.

Coupling the two sides introduces a feedback: intercellular CO$_2$ sets the
assimilation rate, while conductance sets how much CO$_2$ can diffuse in. The
operating point is defined implicitly and must be found at every evaluation
\cite{Damour2010}. Iterative solutions are widely used, and analytic solutions
for particular stomatal formulations are known \cite{Buckley2017,Kim2003}. The
Farquhar and Jarvis equations used here are inherited physiological components,
not a contribution of this paper. The contribution is how they are solved at
runtime and coupled to the carbohydrate, nutrient and biomass states.

\subsection{Diel carbohydrate turnover and photoperiod}

Carbohydrate dynamics link assimilation to growth. In the light, assimilated
carbon is partitioned between mobile sucrose and transient leaf starch; in the
dark, starch degradation supplies carbon for respiration, maintenance and
growth \cite{Stitt2012,Smith2007}. In \textit{Arabidopsis thaliana} this
degradation is closely regulated and the reserve is almost completely used by
dawn \cite{Graf2010,Scialdone2013,Webb2015}. Models of the process connect
circadian regulation, sucrose availability and sugar signalling to partitioning
and growth \cite{Feugier2013,Seki2017}, and the separation between cytosolic and
vacuolar sucrose has been described in detail \cite{Lohaus2000,Nadwodnik2008}.

Photoperiod changes the balance between daytime accumulation and nocturnal
consumption, and \textit{Arabidopsis thaliana} adjusts partitioning and growth
to day length \cite{Sulpice2014,Mengin2017,Annunziata2017}. This makes
photoperiod scenarios a suitable evaluation setting: they perturb both phases of
the diel carbon economy through one variable, with documented responses. Growth
also depends on nutrient status, and nitrogen and phosphorus enter mechanistic
models as internal states or limiting factors
\cite{Gastal2002,Xu2012,Barber1995,Shen2011,Carstensen2018}. Models that
resolve the diel cycle in detail, however, usually treat assimilation as given
\cite{Scialdone2013,Graf2010,Feugier2013}, so the day--night balance and the
gas exchange that produces it are not solved as one system.

\subsection{Verification, validation and sensitivity practice in simulation}

Two distinctions from simulation methodology are relevant here. The first is
between verification and validation: agreement with measurements does not show
that the numerical procedures inside the simulator are correct
\cite{Zeigler2000,Cellier1991}, so residual checks against analytic solutions,
integration refinement and bookkeeping checks address a different question from
data agreement. The second is between calibration residual and predictive
error: when a model is fitted and evaluated against the same small set of
aggregate observations, the reported error is a calibration residual, and
reported accuracy depends strongly on the calibration protocol
\cite{Maiorano2017,Wallach2018,Holzworth2014,Jones2003}. Ablation attributes
behaviour to model structure by removing, simplifying or replacing a component
\cite{Zeigler2000}. Variance-based sensitivity methods give first- and
total-order indices and quantify interaction structure
\cite{Sobol2001,Saltelli2010}; Morris screening indicates nonlinearity and
interaction more cheaply \cite{Morris1991,Campolongo2007}, Latin hypercube
designs are standard for sampling \cite{McKay1979}, and testing the estimator
against a function with closed-form indices is good practice
\cite{Ishigami1990}. These indices describe model behaviour over the declared
parameter domain; they do not establish physiological importance, and screening
measures are not an exact variance decomposition.

\subsection{Open problems addressed here}

Read together, this literature leaves a specific gap. Leaf-level gas-exchange
models solve the supply--demand operating point accurately, but at fixed
carbohydrate status and without downstream dynamics
\cite{Farquhar1980,Jarvis1976,vonCaemmerer2000}. Whole-plant and crop
simulators carry growth over days to seasons, but prescribe or simplify the
carbon input and rarely resolve the diel cycle
\cite{Jones2003,deWit2019,Holzworth2014}. Dedicated diel carbohydrate models
resolve the day--night starch balance, but take assimilation as an input
\cite{Scialdone2013,Graf2010}. Models that do compute assimilation
biochemically and carry it into leaf-area growth are the closest published
relatives \cite{Weraduwage2015}, but they do not carry the diel starch and
sucrose pools as integrated states. The feedback loop connecting conductance,
intercellular CO$_2$, carbohydrate status and growth is therefore opened
somewhere in each family. Two further problems are methodological. The algebraic
constraint created by the coupling must be solved at every step, yet the solver
and its behaviour outside the calibrated range are often left unstated.
Calibration is also frequently performed against the same aggregate observations
used for evaluation, so predictive error cannot be separated from a calibration
residual.

This paper closes part of that gap. We do not propose new process equations for
photosynthesis, stomatal conductance or starch turnover; those are taken from
the literature. We add the coupled formulation with a computed carbon input,
its analytic runtime closure and fast/slow numerical treatment, a verification
protocol independent of data agreement, and an evaluation design in which whole
photoperiods are held out and the fitted parameter set is shared across
conditions. The aim is not to replace crop simulators or
functional--structural models, but to provide a compact and reproducible core
in which gas exchange is propagated explicitly into downstream states.

We also state where the model falls short, with numbers in
Section~\ref{sec:results}. The model does not generate the full observed
photoperiod dependence of growth from its shared dynamics; part of that
dependence remains localised in one growth-conversion coefficient. The diel
starch balance closes only up to an upper photoperiod, above which the reserve
accumulates from day to day (Figure~\ref{fig:envelope}). Canopy geometry, water
transport, temperature acclimation and multi-species parameterisation are
outside the present scope. These are stated as named structural limitations in
Section~\ref{sec:limitations}.

